% =========================================================================
% CHAPTER 1: MỞ ĐẦU VÀ ĐỊNH HƯỚNG NGHIÊN CỨU
% =========================================================================

\section{Mở đầu và Định hướng nghiên cứu}

% --- SLIDE 1.1: Bối cảnh nghiên cứu ---
\begin{frame}{1.1 Bối cảnh nghiên cứu}
    \small
    \begin{columns}[T]
        \column{0.48\textwidth}
        \textbf{Vai trò của mạng Campus:}
        \begin{itemize}
            \item Kết nối người học, giảng viên, phòng ban và các hệ thống dùng chung.
            \item Hỗ trợ học tập, nghiên cứu, quản trị và cung cấp dịch vụ số.
            \item Liên kết nhiều khu vực chức năng và nhiều cơ sở địa lý.
        \end{itemize}

        \column{0.48\textwidth}
        \textbf{Vấn đề đặt ra:}
        \begin{itemize}
            \item Hạ tầng ngày càng lớn, nhiều thiết bị và chính sách truy cập.
            \item Quản trị phân tán làm tăng độ phức tạp khi thay đổi hệ thống.
            \item Kết nối WAN truyền thống thiếu linh hoạt khi mở rộng đa cơ sở.
        \end{itemize}
    \end{columns}
    \vspace{0.2cm}
    \begin{block}{Định hướng nghiên cứu}
        Kết hợp kiến trúc Campus phân cấp với \textbf{SDN} cho mạng nội bộ và \textbf{SD-WAN} cho kết nối đa cơ sở.
    \end{block}
\end{frame}

% --- SLIDE 1.2: Mục tiêu và câu hỏi nghiên cứu ---
\begin{frame}{1.2 Mục tiêu và câu hỏi nghiên cứu}
    \begin{block}{Mục tiêu tổng quát}
        Xây dựng khung thiết kế Campus Network ứng dụng SDN và SD-WAN phù hợp với môi trường đại học đa cơ sở.
    \end{block}
    \vspace{0.15cm}
    \textbf{Các câu hỏi nghiên cứu chính:}
    \begin{enumerate}
        \item Kiến trúc Campus phân cấp tổ chức vai trò của Core, Distribution và Access như thế nào?
        \item SDN có thể tập trung hóa lớp điều khiển và quản lý chính sách trong Campus ra sao?
        \item SD-WAN tổ chức underlay, overlay và kết nối an toàn giữa các cơ sở như thế nào?
        \item Hai công nghệ cần được tích hợp ở vị trí nào trong một mô hình thống nhất?
    \end{enumerate}
\end{frame}

% --- SLIDE 1.3: Đối tượng và phạm vi ---
\begin{frame}{1.3 Đối tượng và phạm vi nghiên cứu}
    \begin{columns}[T]
        \column{0.48\textwidth}
        \textbf{Đối tượng nghiên cứu:}
        \begin{itemize}
            \item Kiến trúc mạng Campus trong môi trường đào tạo.
            \item Phân tầng mạng, VLAN và quy hoạch địa chỉ IP.
            \item Kiến trúc SDN, OpenFlow và SDN Controller.
            \item Kiến trúc SD-WAN, underlay và overlay.
        \end{itemize}

        \column{0.48\textwidth}
        \textbf{Phạm vi của báo cáo tiến độ:}
        \begin{itemize}
            \item Tổng hợp cơ sở lý thuyết liên quan.
            \item Phân tích yêu cầu cho mô hình đại học đa cơ sở.
            \item Đề xuất kiến trúc logic và mô hình mô phỏng.
            \item Ghi nhận các hạng mục đã và đang thực hiện.
        \end{itemize}
        \vspace{0.1cm}
        \textbf{Ngoài phạm vi hiện tại:} triển khai trên hạ tầng vật lý và đo kiểm hiệu năng chuyên sâu.
    \end{columns}
\end{frame}

% --- SLIDE 1.4: Phương pháp nghiên cứu ---
\begin{frame}{1.4 Phương pháp nghiên cứu}
    \begin{enumerate}
        \item \textbf{Nghiên cứu tài liệu:} hệ thống hóa khái niệm Campus Network, SDN, SD-WAN và các giao thức liên quan.
        \vspace{0.12cm}
        \item \textbf{Phân tích yêu cầu:} xác định phạm vi người dùng, khu vực chức năng, dịch vụ và nhu cầu kết nối đa cơ sở.
        \vspace{0.12cm}
        \item \textbf{Thiết kế mô hình:} xây dựng kiến trúc phân tầng, phân vùng logic và quan hệ giữa SDN với SD-WAN.
        \vspace{0.12cm}
        \item \textbf{Mô phỏng triển khai:} dựng mô hình trên EVE-NG, cấu hình từng vùng (LAN/SDN, chi nhánh, WAN) và kiểm chứng hoạt động của các luồng kết nối.
    \end{enumerate}
    \vspace{0.15cm}
    \begin{block}{Trọng tâm giai đoạn hiện tại}
        Triển khai thực nghiệm trên mô hình EVE-NG: vận hành vùng SDN và kết nối đa site SD-WAN trước khi tổng hợp kết quả.
    \end{block}
\end{frame}
